\begin{lemma}[Residual future transport and genuine row restoration]
In a finite disjoint-union ambient set, let
$U\cap K=U\cap Z=K\cap Z=V\cap Z=\varnothing$,
$|U|=|K|=n$, $|V|=|Z|=n+1$, and $n\ge1$.
Let $v_0\in V$ and let $\tau:U\to V\setminus\{v_0\}$ be a bijection
such that $uv\in M\Longleftrightarrow v=\tau(u)$ on $U\times V$.
Assume $W\subseteq O$, $J\subseteq H$, $e_0\le1$,
$O[V,Z]$ is $(e_0,p_0)$-super-regular, and $W[V,Z]$ is
$(\xi,p)$-super-regular. Suppose $p\le1$, $ph\le1$, and
\[
\xi<\eta\le1,\quad\xi+2/n\le\eta,\quad
\xi+1/(\eta n)\le\eta,
\]
\[
\omega<e\le1,\quad\omega+2/n\le e,\quad
\omega+1/(en)\le e,\quad ph+e_0+1/n\le p_0.
\]
There is an $(\eta,p)$-super-regular graph $F[U,Z]$ with
$uz\in F\Longleftrightarrow\tau(u)z\in W$ on $U\times Z$.
For every bijection $\sigma:U\to K$, if the matching-induced graph on
subtype copies of $U,Z$, with edges $uz\in F$ and $\sigma(u)z\in J$,
is $(\omega,ph)$-super-regular, the graph
\[
T_\sigma:\quad vz\in E(T_\sigma)\Longleftrightarrow
v\in V,\ z\in Z,\ vz\in O,\quad
\forall u\in U\ (uv\in M\Rightarrow\sigma(u)z\in H),
\]
symmetrized as a simple graph, contains a spanning
$(e,ph)$-super-regular graph on $V,Z$.
\end{lemma}
\begin{proof}
Before testing any edges, use \noderef{MatchingInducedLeftRowDeletion}
on $W[V,Z]$ to delete the row $v_0$. The size parameter is $n$, the
right part has $n+1$ vertices, the errors are $\xi,\eta$, and all the
unit and slack hypotheses are displayed. The resulting pair on
$V\setminus\{v_0\},Z$ is $(\eta,p)$-super-regular.
Pull it back along $\tau$ and restrict to $U,Z$. Bijections preserve
all the degrees and tested rectangle edge counts in
\noderef{SuperRegularPair}, giving the asserted graph $F$ and its
exact cross-part adjacency equivalence.

For a successful $\sigma$, define the residual graph $B$ to have just
the edges $\tau(u)z$ with $uz\in F$ and $\sigma(u)z\in J$, and their
reverses. Apply \noderef{MatchingInducedUpdatedCandidacyTransport} to
$\tau$, $\sigma$, and $B$, using the complete bipartite matching graph
on $U,K$. Bijectivity makes $\sigma$ a perfect matching as in
\noderef{PerfectMatching}. Its adjacency hypothesis is the definition
of $B$; its exact-image hypotheses are the equivalence for
$V\setminus\{v_0\}=\tau(U)$ and the identity equivalence for $Z$.
Thus $B[V\setminus\{v_0\},Z]$ is $(\omega,ph)$-super-regular.
The pulled-back-row identity, $W\subseteq O$, and $J\subseteq H$
show $B\subseteq T_\sigma$. The source-matching equivalence ensures
there is precisely one newly imposed host constraint on every residual
row. It also ensures that $v_0$ has no neighbour in $U$, so its entire
old $O$-row is still available in $T_\sigma$.

Put $k=\lfloor ph|Z|\rfloor$. The outer degree reserve implies
$1/n\le e$ and hence $e|Z|\ge1$, since $|Z|=n+1$.
Integer rounding therefore gives
$(ph-e)|Z|\le k\le(ph+e)|Z|$.
The old degree lower bound gives
$\deg_{T_\sigma}(v_0,Z)=\deg_O(v_0,Z)\ge(p_0-e_0)|Z|$.
The gap inequality implies $p_0-e_0\ge ph+1/n$, and consequently
$k\le ph|Z|\le(p_0-e_0)|Z|\le\deg_{T_\sigma}(v_0,Z)$.
Here $e_0$ is finite because $e_0\le1$; $p_0$ is finite because otherwise
the old row lower bound at $v_0$ would be infinite, contradicting its
finite degree and $|Z|>0$. Thus subtracting $e_0$ in the gap inequality
is legitimate. These calculations also cover zero output density.
Now apply \noderef{FutureExtraVertexStability} to $B$ inside
$T_\sigma$, on residual left size $n$ and right size $n+1$, with
row size $k$, errors $\omega,e$, and density $ph$. Its two analytic
reserves and unit bounds are exactly the hypotheses. It selects a
$k$-subset of the genuine old row, adjoins that row, and returns a
spanning contained $(e,ph)$-super-regular graph. No artificial row and
no surjection onto the full future part are used in this case.
\end{proof}
